\begin{claim}{AC.36}{Proved: an exact nonabelian bounded-exponent example}
Let \(L=\mathbf F_{4096}\), let \(D/L\) be the Fermat curve
\[
                         X^{39}+Y^{39}+Z^{39}=0,
\]
and choose a primitive thirteenth root \(\zeta\in L\).
The automorphisms
\[
 \sigma(X:Y:Z)=(\zeta X:\zeta^3Y:\zeta^9Z),
 \qquad \tau(X:Y:Z)=(Y:Z:X)
\]
generate a constant group \(G\simeq C_{13}\rtimes C_3\)
acting freely on \(D\). For \(C=D/G\), one has
\[
                  g(D)=703,\qquad g(C)=19.
\]
For every odd positive integer \(m\),
\[
 \begin{aligned}
 J_D(\mathbf F_{4096^m})[2^\infty]
                  &\simeq(\mathbf Z/2)^{24},\\
 J_C(\mathbf F_{4096^m})[2^\infty]
                  &\simeq(\mathbf Z/2)^4.
 \end{aligned}
\]
Thus this cover has relative class-number valuation
\[
             v_2\!\left(\frac{h(D)}{h(C)}\right)=20
\]
over every such odd extension: it fails (R).
Nevertheless every \(P\in C(\mathbf F_{4096^m})\), with
\(m\) odd, has a rational lift to this same cover after
an additional extension of degree in \(\{1,3,13\}\),
and the bounded-exponent nonabelian criterion applies there,
since \(39>(2-1)^2\).

For an even extension degree \(m\), the exact exponent of
\(J_D(\mathbf F_{4096^m})[2^\infty]\) is
\(2^{v_2(m)+2}\). In particular it is eight for \(m=2\);
the numerical inequality for this order-39 cover then fails.
The example distinguishes two sufficient tests on a specified
cover; it does not assert that \(C\) has no other (R) witness.
\end{claim}
\begin{proof}
The relation is \(\tau\sigma\tau^{-1}=\sigma^3\).
Every nonidentity element of \(\langle\sigma\rangle\)
has three distinct projective eigenvalues, and its fixed
points are the coordinate vertices, none on \(D\).
Every remaining nonidentity element of \(G\) has order
three and is conjugate to \(\tau\) or \(\tau^2\).
Their projective fixed points have coordinates
\((1:\omega:\omega^2)\), with \(\omega^3=1\),
and at such a point the Fermat expression is
\(1+1+1=1\). Thus the action is free.
The plane genus formula and the etale Riemann--Hurwitz
formula give the two stated genera.

Here are the complete spectral data, so that the claim
about the integral two-primary exponent is independently
checkable. Put \(q=4096\), choose a multiplicative
character \(\chi:L^\times\to\mu_{39}\), and set
\[
 J(a,b)=\sum_{x\in L\setminus\{0,1\}}
                   \chi(x)^a\chi(1-x)^b,
 \qquad \alpha_{a,b}=-J(a,b).
\]
The indices are \(1\le a,b\le38\), \(a+b\not\equiv0\pmod{39}\).
For the diagonal group \(\mu_{39}^2\), each such character
occurs exactly once in \(H^1\), and its Frobenius scalar is
\(\alpha_{a,b}\). This assertion, including its sign and its
validity for both prime-to-two etale and crystalline cohomology,
is exactly \cite[Corollary~2.4, p.~181]{AC_Katz1981}.
The integral crystalline version is explained on p.~183;
Proposition~4.2 on pp.~185--186 also transfers the character
polynomial identities to the Jacobian endomorphism ring.

The accompanying program evaluates all \(1406\) sums in
\(\mathbf Z[\zeta_{39}]\), using the binary field
\[
              L=\mathbf F_2[t]/(t^{12}+t^3+1)
\]
and the primitive multiplicative generator encoded by the
integer \(3\). It checks each equality
\(\alpha\bar\alpha=4096\), each invariance under
\(\zeta_{39}\mapsto\zeta_{39}^2\), and the following table.
The conductor is \(39/\gcd(a,b,39)\); the second column
is \(\alpha+\bar\alpha\), including twice the scalar in
the rational cases.
\[
\begin{array}{r|r|r}
 \text{conductor}&\alpha+\bar\alpha&\text{character pairs}\\ \hline
 3&128&2\\
 13&-128&132\\
 39&-128&360\\
 39&-80&216\\
 39&-28&72\\
 39&28&432\\
 39&80&144\\
 39&115&48
\end{array}
\]
In particular the possible trace \(-115\) does \emph{not}
occur. As an independent sign check, direct enumeration
of projective points gives \(\#D(L)=30537\), and the
sum of the displayed eigenvalues is \(-26440=q+1-30537\).
Consequently the zeta numerator is
\[
\begin{aligned}
 P_D(T)={}&(1-64T)^2(1+64T)^{492}\\
 &\cdot(1+80T+4096T^2)^{108}
        (1+28T+4096T^2)^{36}\\
 &\cdot(1-28T+4096T^2)^{216}
        (1-80T+4096T^2)^{72}\\
 &\cdot(1-115T+4096T^2)^{24}.
\end{aligned}
\]

It is essential to prove an exponent bound on the actual
Tate lattice, rather than infer one from an unspecified
isogeny decomposition. Let \(F\) denote the \(L\)-Frobenius
endomorphism of \(J_D\). It is central in the semisimple
algebra \(\operatorname{End}_L^0(J_D)\); hence its minimal
polynomial over \(\mathbf Q\) is square-free. By the
faithful prime-to-two Tate realization, its roots occur
among the reciprocal roots of the displayed zeta numerator. The action
of \(F\) on \(V_2^{\rm et}(J_D)\) remains semisimple,
and all its eigenvalues are two-adic units, since \(F\)
is an automorphism of its integral etale Tate lattice.

There is only one unit root among the corresponding
Frobenius polynomials: the root \(\beta\in\mathbf Z_2^\times\) of
\[
                       X^2-115X+4096=0.
\]
The other root \(\gamma\) has valuation twelve.
For the traces \(\pm80\) and \(\pm28\), both roots
have positive valuation, as do the rational roots
\(\pm64\). Thus semisimplicity gives
\[
               F=\beta\,\operatorname{id}
               \quad\text{on }V_2^{\rm et}(J_D),
\]
and therefore also on the actual integral Tate lattice.
This argument excludes any hidden lattice extension.
Since
\[
 (1-\beta)(1-\gamma)=4097-115=3982=2\cdot1991,
\]
one has \(v_2(\beta-1)=1\). The rational two-primary
subgroup is therefore elementary abelian. Its order is
\(2^{24}\), directly from \(P_D(1)\), proving its exact
structure without needing to infer it merely from
Newton slopes.

For completeness, the quotient calculation is also made
at the level of individual characters. Write
\(c=-a-b\pmod{39}\). The \(\sigma\)-invariant characters
are exactly those with
\[
                    a+3b+9c\equiv5a+7b\equiv0\pmod{13}.
\]
There are \(110\) of them. Under the cycle
\((a,b,c)\mapsto(b,c,a)\), two are fixed,
namely \((13,13,13)\) and \((26,26,26)\), and the
other \(108\) form \(36\) orbits of length three.
Each length-three orbit contributes one invariant line,
with its common Jacobi eigenvalue. The two fixed lines
factor through the Fermat cubic. The coordinate cycle
acts freely on that genus-one curve and hence by a
translation, so it acts trivially on these two lines.
The invariant eigenvalues are \(64\) twice, \(-64\)
twenty-eight times, and the two roots of
\(X^2-115X+4096\) four times each. Thus
\[
        P_C(T)=(1-64T)^2(1+64T)^{28}
                           (1-115T+4096T^2)^4.
\]
The same scalar argument, or the pullback embedding
on two-adic Tate modules for the odd-degree cover,
gives the elementary abelian group of order \(2^4\).
In particular \(f_2(D)=24\), \(f_2(C)=4\), and
\(v_2(P_D(1)/P_C(1))=24-4=20\).

The field claims now follow with no change of cover.
The root \(\beta\) is congruent to \(115\pmod{4096}\),
since \(\beta=115-\gamma\). Therefore
\[
                   v_2(\beta-1)=1,\qquad
                   v_2(\beta+1)=2.
\]
For odd \(m\), the factor
\(1+\beta+\cdots+\beta^{m-1}\) is a unit, so
\(v_2(\beta^m-1)=1\). For even \(m\), the elementary
two-adic lifting formula gives
\[
                  v_2(\beta^m-1)=v_2(m)+2.
\]
As \(F\) is scalar on the lattice, these are exact
exponents of the finite rational groups, not only bounds.

For any rational point of \(C\) over an odd-degree
extension of \(L\), its geometric fibre is a \(G\)-torsor.
Frobenius acts there as multiplication by some element
of \(G\), up to conjugacy; its order is one, three or
thirteen. Thus an additional extension of that odd
degree makes the fibre split. The exponent stays two.
The nonabelian annihilator argument of AG.5, with
\(a=2\cdot1\in\mathbf Z[G]\), now excludes a collision
because its degree condition is \(39>(2-1)^2\).
For even extension degree two the exponent is eight
and \(39\le(8-1)^2=49\), which explains the stated
limitation of this particular numerical test.
\end{proof}

\paragraph{Scope of the example.}
There is a separate geometric reason that this particular
quotient does not provide a new unresolved instance of A1.
The two power maps from the Fermat curve of degree \(39\)
to the Fermat curves of degrees three and thirteen are
jointly injective on geometric points: in any common
nonzero coordinate chart, equality of both images forces
ratios whose third and thirteenth powers are both one,
hence forces equal ratios. Both target Jacobians have
two-rank zero, as is already shown by the conductor-three
and conductor-thirteen entries of the Jacobi table.
A two-primary point difference on \(D\) therefore maps
to zero in both Jacobians; the Abel--Jacobi embeddings
and joint injectivity force the two points to coincide.
The collision-lifting statement AC.1 transfers this
exclusion to \(C\), over every finite extension. The
point of AC.36 is the exact nonabelian exponent test,
including a cover for which (R) fails; it is not a claim
that all-extension injectivity of this example required
that numerical test.
